\section{The far exterior at growing height}
\label{sec:gt-height}

The finite lower-mass bounds use information about individual groups.
At growing height a cumulative count can pay the contribution of the
far groups together. We use physical time \(t=\lambda\), and write
\(M=-K\), so the physical gap law is \(D'=2-4MD\). For a pair with
real coordinates \(c=(z_1+z_2)/2\) and \(D=(z_1-z_2)^2/4\),
Proposition~\ref{gt-exterior-proposition} expresses \(M\) through the
complete exterior divisor, after removal of the two target occurrences.

\subsection{Uniform counts, including time zero}

The strip bound \eqref{gt-canonical-strip} holds on the full physical
slab \(0\le t\le1/2\) by the same de Bruijn theorem. The order bound
of Section~\ref{gt-source-curvature} also holds uniformly on this
slab, since \(u^2/2\) in place of \(u^2/10\) is still absorbed by
the double-exponential decay in \eqref{eq:tail}. The source integral
has two further consequences. First, \(H_t(\ii y)>0\) for real
\(y\), by positivity of the integrand \(W_t(u)\cosh(yu)\) on the
half-line. Second, \(H_t\to H_0\) locally uniformly in \(z\) as
\(t\downarrow0\). On compact \(z\)-sets the cosine factor is bounded
by \(\e^{Bu}\) for a fixed \(B\), so the majorant
\(\Phi(u)\e^{u^2/2+Bu}\) gives this convergence.

Let \(\nu_t=\sum_\varrho\delta_{\Re\varrho}\) be the projection of the
complete divisor of \(H_t\), with each occurrence counted, and put
\(N_t(X)=\nu_t((0,X])\). For \(X\ge4\pi\), define
\begin{equation}
 G_t(X)=\frac{X}{4\pi}\log\frac{X}{4\pi}-\frac{X}{4\pi}
        +\frac{11}{8}+\frac{t}{16}\log\frac{X}{4\pi},
 \qquad E_t(X)=N_t(X)-G_t(X).
 \label{gt-height-main-count}
\end{equation}
For \(0<t\le1/2\), the normalization in \eqref{eq:field} gives the
count of Theorem~1.5(iv) of Polymath, arXiv v2, pp.~9--10, with proof
on p.~61. It is a multiplicity count of all complex zeros with
positive real part at most \(X\), with an absolute discrepancy
bound \(C_0\log(2+X)\), allowing an absolute enlargement for the
right-continuous convention. Part~(iv) applies throughout this range
of \(X,t\). \cite[Theorem~1.5(iv)]{polymath}

\begin{lemma}[Complete projected counts on the closed slab]
\label{gt-height-count}
There is an absolute constant \(C_0\) such that
\begin{equation}
 |E_t(X)|\le C_0\log(2+X)\qquad (X\ge4\pi,\ 0\le t\le1/2).
 \label{gt-height-discrepancy}
\end{equation}
For an absolute \(C_{\rm u}\ge1\) one also has
\begin{equation}
 \nu_t([U,U+1])\le C_{\rm u}\log(2+|U|)
 \qquad(U\in\R,\ 0\le t\le1/2).
 \label{gt-height-unit-count}
\end{equation}
In particular \(N_t(X)\ll X\log(2+X)\) for \(X\ge4\pi\).
\end{lemma}
\begin{proof}
Fix \(X>4\pi\) and choose a sequence \(\epsilon\downarrow0\) such that
\(X-\epsilon>4\pi\) and \(X\pm\epsilon\) avoid the real coordinates
of the \(H_0\) zeros. Local finiteness follows from the strip and
isolated zeros. At each chosen endpoint \(Y=X\pm\epsilon\), take
the rectangle with vertical sides at \(0,Y\) and horizontal sides
at \(\Im z=\pm2\). The boundary is zero-free for \(H_0\). Local
uniform convergence and Rouch\'e's theorem give the same count
inside for all sufficiently small positive \(t\) at that \(Y\).
By the strip and the zero-free vertical sides, this is equality
of \(N_t(Y)\) and \(N_0(Y)\). Letting \(t\downarrow0\) at the two
fixed endpoints uses \(G_t(Y)\to G_0(Y)\). Monotonicity then gives
\[
 \begin{aligned}
 G_0(X-\epsilon)-C_0\log(2+X-\epsilon)
 &\le N_0(X-\epsilon)\le N_0(X),\\
 N_0(X)&\le N_0(X+\epsilon)
    \le G_0(X+\epsilon)+C_0\log(2+X+\epsilon).
 \end{aligned}
\]
Taking \(\epsilon\downarrow0\) bounds the full cumulative count even
when \(X\) carries an atom. At \(X=4\pi\) use a bracket from above
for the upper bound and \(N_0(4\pi)\ge0\) for the lower bound,
enlarging the absolute constant.

For \(X>4\pi\), the left limit \(E_t(X-)\) inherits
\eqref{gt-height-discrepancy}, so the atom
\(J_t(X)=\nu_t\{X\}=E_t(X)-E_t(X-)\) is \(O(\log(2+X))\).
At \(X=4\pi\) it is bounded by \(N_t(4\pi)=O(1)\).
For \(U>4\pi\), subtract \(N_t(U-)\) from \(N_t(U+1)\) and use
the bound on the increment of \(G_t\), given by
\[
 G_t'(X)=\frac{1}{4\pi}\log\frac{X}{4\pi}+\frac{t}{16X}.
\]
For \(0\le U\le4\pi\) use \(N_t(4\pi+1)=O(1)\) and the absence
of zeros with zero real part. Evenness gives the bound on negative
intervals by reflecting all occurrences, and intervals straddling
zero are in the bounded range. This proves \eqref{gt-height-unit-count},
including both endpoints of each unit interval. The asserted upper
bound on \(N_t\) follows directly from \eqref{gt-height-main-count}
and \eqref{gt-height-discrepancy}.
\end{proof}

\subsection{The signed far-tail main term}

\begin{theorem}[Uniform far exterior of a counted pair]
\label{gt-height-far}
At a slice of the canonical flow with \(0\le t\le1/2\), suppose a
conjugation-invariant contour counts exactly the target pair and
has a separately established zero-free collar. Assume \(c\ge16\pi\),
\(|D|\le1\), and \(3\le R\le c/2\). For \(\varrho=c+x+\ii y\),
use \(q(x,y,D)\) from \eqref{gt-conjugate-weight} and define
\begin{equation}
 S_R=\sum_{\varrho\ {\rm exterior},\,|x|\le R}q(x,y,D),
 \qquad T_R=M-S_R .
 \label{gt-height-signed-split}
\end{equation}
Each sum over zeros counts occurrences. All far weights are
positive, and with \(B_c=\log(2+c)\),
\begin{equation}
 \left|T_R-\frac{\log(c/(4\pi))}{2\pi R}\right|
 \le C\left(\frac{B_c}{R^2}+\frac{B_c}{R^3}+\frac{B_c}{c}\right).
 \label{gt-height-tail-estimate}
\end{equation}
The constant is absolute and unevaluated. The estimate applies to
positive, negative and zero \(D\). \cite{gtfar}
\end{theorem}
\begin{proof}
In \eqref{gt-exterior-k} the reciprocal product has real part
\(q(x,y,D)\). The sum is absolutely convergent and conjugation
invariant, so a nonreal conjugate pair contributes two copies of
this weight and a real occurrence contributes one. The same formula
holds at a double. The roots satisfy
\(|\Re z_j-c|\le1\), since \((z_j-c)^2=D\), so the strict far cut
\(|x|>R\) already omits the target. Put \(h=y^2+D\). For
\(|x|\ge3\) one has \(|h|\le2\) and \(x^2-h\ge(7/9)x^2\).
Thus \(q>0\) there, and
\begin{equation}
 q(x,y,D)-x^{-2}
 =\frac{(D-3y^2)x^2-h^2}{x^2\bigl((x^2-h)^2+4x^2y^2\bigr)}
 =O(|x|^{-4}).
 \label{gt-height-tail-kernel}
\end{equation}
The fourth-power error comes from taking the real part before
approximation. Its numerator is \(O(x^2)\) and its denominator
is at least a positive constant times \(x^6\).

Write \(u=\Re\varrho\), \(\alpha=c/2\) and \(\beta=3c/2\).
Each of the two parts of a shell
\(R+j<|u-c|\le R+j+1\) intersected with \([\alpha,\beta]\)
fits inside a closed unit interval with start of size \(O(c)\).
The bound \eqref{gt-height-unit-count} therefore gives
\[
 \sum_{\substack{u\in[\alpha,\beta]\\ |u-c|>R}}
     |u-c|^{-4} \ll B_c\sum_{j\ge0}(R+j)^{-4}
 \ll B_c/R^3 .
\]
For \(p\in\{2,4\}\), the range \(0<u<\alpha\) costs
\(O(B_c/c^{p-1})\) for the inverse-power sum, by
\(|u-c|\ge c/2\) and \(N_t(\alpha)\ll cB_c\).
For \(u>\beta\) one has \(u-c\ge u/3\), and the dyadic range
\(2^j\beta<u\le2^{j+1}\beta\) contains
\(O(2^jc(B_c+j))\) occurrences. Its weighted bounds sum to
\(O(B_c/c^{p-1})\), since \(\sum_{j\ge0}2^{-(p-1)j}(1+j)<\infty\).
For \(u=-s<0\), the weight is \((c+s)^{-p}\). By reflection the
range \(0<s\le c\) has \(O(cB_c)\) occurrences, and dyadic ranges
in \(s>c\) give the same bound as on the upper positive side.
This includes the reflected target and all other negative-coordinate
zeros. There is no occurrence with \(u=0\). Thus
\eqref{gt-height-tail-kernel} and these bounds give
\begin{equation}
 T_R=\int_{[\alpha,\beta]\cap\{|u-c|>R\}}
         \frac{\dd\nu_t(u)}{(u-c)^2}
       +O(B_c/R^3+B_c/c).
 \label{gt-height-tail-central}
\end{equation}

For \(R<c/2\), set \(a=c-R\), \(b=c+R\) and
\(f(u)=(u-c)^{-2}\). The central domains are
\([\alpha,a)\) and \((b,\beta]\).
On these intervals \(\dd E_t=\dd\nu_t-G_t'(u)\,\dd u\).
Stieltjes integration by parts with the specified endpoints gives
\begin{align*}
 \int_{[\alpha,a)} f\,\dd E_t
  &= f(a)E_t(a-)-f(\alpha)E_t(\alpha-)
      -\int_\alpha^a f'(u)E_t(u)\,\dd u,\\
 \int_{(b,\beta]} f\,\dd E_t
  &= f(\beta)E_t(\beta)-f(b)E_t(b)
      -\int_b^\beta f'(u)E_t(u)\,\dd u .
\end{align*}
These formulas include atoms at \(\alpha,\beta\) and exclude
atoms at \(a,b\), which belong to \(S_R\). Each excluded atom has
inverse-square weight \(R^{-2}J_t(a)\) or \(R^{-2}J_t(b)\), both
\(O(B_c/R^2)\). In the formulas all discrepancy values are
\(O(B_c)\), the endpoint weights are at most \(R^{-2}\), and
the variation of \(f\) on either interval is at most \(R^{-2}\).
The cost of replacing \(\dd\nu_t\) by \(G_t'(u)\,\dd u\)
in \eqref{gt-height-tail-central} is therefore \(O(B_c/R^2)\).
At \(R=c/2\) both central intervals are empty and the coincident
endpoints belong to \(S_R\).

On the central band the derivative formula in the count lemma gives
\[
 G_t'(u)=\frac{\log(c/(4\pi))}{4\pi}
          +O\!\left(\frac{|u-c|+1}{c}\right)
\]
by the mean-value theorem and \(u\ge c/2\). After weighting,
the error is \(O(c^{-1}\int_R^{c/2}(s+1)s^{-2}\,\dd s)
=O(B_c/c)\). The constant density contributes
\[
 \frac{\log(c/(4\pi))}{2\pi}\int_R^{c/2}\frac{\dd s}{s^2}
 =\frac{\log(c/(4\pi))}{2\pi R}-\frac{\log(c/(4\pi))}{\pi c}.
\]
The final term is \(O(B_c/c)\), which proves
\eqref{gt-height-tail-estimate} also at \(R=c/2\).
\end{proof}

For a fixed sufficiently large physical \(R\),
Theorem~\ref{gt-height-far} gives \(T_R\asymp(\log c)/R\) as
\(c\to\infty\). Relative to \((\log c)/R\) the error is
\(O(R^{-1}+R^{-2}+R/c)\), so choosing \(R\) first and then
increasing \(c\) gives the lower and upper bounds. For a frozen
reference height \(T\), put \(L=\log T\), \(\ell=1/L\) and
\(R=\sqrt L\). If \(T/2\le c\le2T\) and \(L\ge9\), the cutoff
is admissible and \(\log(c/(4\pi))=L+O(1)\). Consequently
\begin{equation}
 T_R=\frac{\sqrt L}{2\pi}+O(1),\qquad
 \ell^2 T_R=\frac{1}{2\pi L^{3/2}}+O(L^{-2}).
 \label{gt-height-normalized-tail}
\end{equation}
The positive far interaction grows in physical units and tends to
zero after normalization. The constants give an asymptotic range
with unevaluated onset. As \(c\) moves, the two sharp-cutoff sums
can change membership, exchanging equal and opposite weights.
At a slice the total weight on the two cutoff lines is \(O(B_c/R^2)\)
by \eqref{gt-height-tail-kernel} and the atom bound. The gap law
uses the original regular sum \(M\).

In the backward cost \eqref{eq:gt-backward-gap} the exponent is
\(-4\int_t^v M(s)\,\dd s\). The far estimate bounds one component
of this integral. To infer a positive gap on a new tube one still
needs an enclosure of the signed \(S_R\) sufficient for the whole
cost, and an independently obtained strict incoming margin at \(v\).
Direct deflation on that tube uses the same count, collar and
pair coordinates, the same signed near formula and far estimate,
and the same near and incoming data when acquired.

\subsection{Identical horizontal counts and opposite signs}

The kernel \eqref{gt-conjugate-weight} changes sign with the
imaginary offset of an exterior zero close to the center.
The following fixed-slice example isolates the information lost
by projecting the divisor to the real coordinate. It uses the
time-zero counts from Lemma~\ref{gt-height-count}.

\begin{proposition}[An exact-count obstruction at a double]
\label{gt-height-exact-count}
For each \(c\ge\e^4\), put \(\ell=1/\log c\) and \(r=2\ell\).
There are two real even entire functions of order at most one and
zero strip \(|\Im z|\le1\), with exactly the same complete projected
counting measure. Both satisfy the time-zero cumulative law and
closed unit upper counts with constants uniform in \(c\). They
have an isolated double at \(c\), count two in \(|z-c|<\ell\)
and a common zero-free collar \(\ell\le|z-c|\le3\ell/2\).
For their complete deflated interactions \(M_+(c),M_-(c)\),
\begin{equation}
 \ell^2 M_+(c)=\frac12+O((\log c)^{-1}),\qquad
 \ell^2 M_-(c)=-\frac1{16}+O((\log c)^{-1}).
 \label{gt-height-count-signs}
\end{equation}
For \(r\le R\le3\) the normalized near sums equal \(1/2\) and
\(-1/16\) exactly. \cite{gtcounts}
\end{proposition}
\begin{proof}
Erase the full multiset \(\mathcal E_c\) of \(H_0\) zeros in
the two closed bands \(|\Re z-c|\le3\) and \(|\Re z+c|\le3\).
These disjoint bands contain \(O(\log c)\) occurrences by
\eqref{gt-height-unit-count}, and the multiset is invariant
under negation and conjugation. The quotient
\[
 B(z)=H_0(z)\Big/\prod_{\varrho\in\mathcal E_c}(z-\varrho)
\]
extends through every apparent pole with nonzero value because
all multiplicities there have been divided out. It is real even,
with exactly the remaining divisor. The quotient has order at
most one because outside a sufficiently large disk, at fixed
\(c\), each denominator factor has modulus at least one.
Form the replacements
\begin{align*}
 F_+(z)&=B(z)(z^2-c^2)^2(z^2-(c+r)^2)^2,\\
 F_-(z)&=B(z)(z^2-c^2)^2
          \bigl((z-(c+r))^2+3r^2\bigr)
          \bigl((z+(c+r))^2+3r^2\bigr).
\end{align*}
Polynomial multiplication preserves order at most one. Both
functions are real even. Besides the doubles at \(\pm c\),
\(F_+\) inserts real doubles at \(\pm(c+r)\), and \(F_-\)
inserts conjugate pairs with the same real coordinates and
imaginary coordinates \(\pm\sqrt3 r\). All the insertions
lie in the strip, since \(r\le1/2\).
The nearest exterior insertion to \(c\) is at distance \(r\)
for \(F_+\) and \(2r\) for \(F_-\). Retained zeros have horizontal
distance greater than three and the reflected insertions have
distance at least \(2c\). Both disks and closed collars
therefore have the stated geometry.

The projected counting measures are exactly
\[
 \nu_\pm=\nu_0-\sum_{\varrho\in\mathcal E_c}\delta_{\Re\varrho}
      +2\bigl(\delta_c+\delta_{-c}+\delta_{c+r}+\delta_{-(c+r)}\bigr).
\]
They agree on all Borel sets, including counts with either endpoint
convention. Each unit interval has at most eight added occurrences,
so \eqref{gt-height-unit-count} persists after absolute enlargement.
For the cumulative positive count the correction is zero for
\(X<c-3\), and otherwise is at most \(\#\mathcal E_c+8
=O(\log c)=O(\log(2+X))\). Thus
\(N_\pm(X)=G_0(X)+O(\log(2+X))\) for \(X\ge4\pi\), with uniform
constants. The bound also holds on \(0\le X<4\pi\) by the unit
counts and the continuous extension of \(G_0\) there.

At \(D=0\), Proposition~\ref{gt-exterior-proposition} gives
\(M_\pm(c)=\sum_{\varrho\ {\rm exterior}}(c-\varrho)^{-2}\).
For the common retained divisor,
\[
 \sum_{\varrho\in Z(B)}|c-\varrho|^{-2}
 \le 2C_{\rm u}\sum_{j=3}^\infty
         \frac{\log(2+c)+\log(j+2)}{j^2}=O(\log c).
\]
Here the horizontal shell \(j\le|\Re\varrho-c|<j+1\) fits
inside two closed unit intervals with starts of magnitude at
most \(c+j+1\). The bound includes all retained negative-coordinate
zeros. The double at \(-c\) contributes \(1/(2c^2)\), and the
two occurrences with real part \(-(c+r)\) contribute at most
\(2(2c+r)^{-2}\) in absolute value. The exterior insertions
with real part \(c+r\) contribute \(2/r^2\) for \(F_+\) and
\[
 \frac{1}{(r+\ii\sqrt3 r)^2}+\frac{1}{(r-\ii\sqrt3 r)^2}
       =-\frac{1}{4r^2}
\]
for \(F_-\). Multiplying by \(\ell^2=r^2/4\) gives
\eqref{gt-height-count-signs}. Those two occurrences are
the entire signed near sum for \(r\le R\le3\). For a larger
cutoff the added subset of the retained divisor still costs
at most \(O(\log c)\), and the reflected insertions at most
\(O(c^{-2})\).
\end{proof}

In this fixed-slice analytic-divisor class the common count, strip,
order and collar data cannot determine the sign of \(M\), or give
an \(o(1)\) upper bound on the negative part of \(\ell^2S_R\)
for the cutoffs in the proposition. Along unbounded \(c\) with
\(H_0(c)\ne0\), each replacement differs from the unchanged theta
field because it vanishes at \(c\). Positive Fourier-source
representation and a common physical heat-time tube remain
untested for these replacements. Source-specific, time-correlated
and all-real-exterior estimates use additional information.
With the complete fields, a direct evaluator has at the double
\[
 -(\log(F_\pm/(z-c)^2))''(c)
  =\left(\frac{F_\pm'''(c)}{3F_\pm''(c)}\right)^2
    -\frac{F_\pm''''(c)}{6F_\pm''(c)}
  = M_\pm(c).
\]
The first equality follows by the Taylor series at \(c\).
The jets distinguish the signs, and the shared projected data
give the same ambiguity for direct evaluation and gap transport.

For comparison, the finite polynomial controls and the finite-atom
or compact finite-profile controls fail the canonical cumulative
growth. In the compact-source cases exponential type gives a
Jensen upper count \(O(X)\) inside \(|z|\le X+1\), which bounds
the cumulative count under the strip premise. A nonconstant
Euler-factor translate with multiplier \(1+\gamma\cos(az)\),
\(0<\gamma<1\), \(a>0\), adds \(aX/\pi+O(1)\) occurrences
at time zero. Its two roots per period \(2\pi/a\) have real
coordinates \((2k+1)\pi/a\) and imaginary coordinates
\(\pm\arcosh(1/\gamma)/a\), with multiplicities added at any
common root of the product. This fails the logarithmic discrepancy
even if the added width is at most one. The quadratic-exponential
controls of order two and the reflected super-Gaussian controls
fail the order premise. The documented endpoint-defective weighted-square,
finite-source, finite-score and finite-prime changes fail the
global strip as in Section~\ref{gt-source-curvature}. The damped
smooth-shape and positive-primitive controls have untested strip
and count premises. The normalized merger retains its RH hypothesis
and untested cumulative count. Reflection rigidity identifies the
unchanged weights within its polynomially weighted square-comb class.
Varied observation kernels fail the literal cosine/Newman field
equation, and fractional-profile measures lack an integer divisor.
The exact-count controls above meet the count and collar premises
of their fixed-slice class. The independent incoming condition in
\eqref{eq:gt-backward-gap} remains necessary for a strict backward
passage, as the even quartic illustrates despite positive exterior
mass. \cite{gtjoint,gtfar,gtcounts}
